\section{Related Work}
\label{rw_section}

Accurately measuring software energy consumption has been a longstanding research challenge, leading to the development of various tools that differ in methodology, level of abstraction and target environments.

Several approaches provide language-specific libraries for energy measurement. For instance, pyRAPL \cite{pyrapl} exposes Intel's RAPL counters through a simple Python interface, while jRAPL \cite{jRAPL} enables energy profiling for Java programs. These tools allow fine-grained energy monitoring directly within the program but require source code modifications and are limited to their respective programming ecosystems. In contrast, GreenMeasurer is language-agnostic and operates externally to the target program, avoiding the need for instrumentation.

Other tools aim for higher-level abstractions, often estimating the environmental impact of software execution. CodeCarbon \cite{benoit_courty_2024_11171501}, for example, estimates $CO_{2}$ emissions and energy consumption across hardware components such as CPU, RAM and GPU. While useful for assessing carbon footprints, such estimation-based tools rely on model approximations when low-level measurement interfaces are unavailable, which can compromise precision.

Recent tools such as Scaphandre \cite{Scaphandre} and Kepler \cite{10254956} address software energy monitoring in modern and containerized environments. Scaphandre collects energy data from Intel RAPL interfaces and exports metrics through Prometheus, enabling system-wide energy monitoring with minimal configuration. Kepler integrates with Kubernetes to estimate pod-level energy consumption based on performance counters and hardware telemetry, providing visibility into energy usage across containers and workloads.

GreenMeasurer, on the other hand, focuses on repeatable measurements of a command, without instrumentation, on bare-metal machines. It reads hardware counters directly and automates temperature calibration, repeated executions and (on Intel CPUs) power capping, which has been used in prior studies. Although it does not support containerized or virtualized systems, this design enables simpler setup, making it particularly suitable for experimental and comparative studies of software energy consumption.

We compare GreenMeasurer with two tools: CodeCarbon \cite{benoit_courty_2024_11171501}, a Python library that is integrated into the program under test, and EnergiBridge, a cross-platform command-line tool that measures the energy of a command from hardware counters. \textbf{[TODO: confirm the description, supported platforms and add the reference of EnergiBridge.]} GreenMeasurer's positioning rests on the following properties, each with its justification:

\begin{itemize}
    \item \textbf{Ease of use.} A command and a few parameters are enough; no code change is needed and the same workflow applies to any language.
    \item \textbf{Ease of extension.} A new CPU or GPU is supported by adding one backend directory (Section~\ref{sec:adding_backend}); the core, the build and the CSV schema adapt automatically.
    \item \textbf{Graceful degradation.} CPU and GPU backends are selected independently; an unsupported one yields \texttt{NOT\_SUPPORTED} instead of an abort or a silent estimate.
    \item \textbf{Testable hardware support.} The backend interface allows fake backends, so the selection and fallback logic is verified by an automated test suite without needing every hardware combination (these tests check orchestration, not counter accuracy).
    \item \textbf{Self-describing output.} Each backend declares its metrics (name, unit, kind, scope) and the CSV header is generated from them, so the output cannot drift from what the hardware provides. A missing value is never written as zero.
    \item \textbf{Controlled conditions.} Temperature calibration and the disabling of Wi-Fi during runs address known sources of variation in energy measurements. Whether they reduce variance is to be shown experimentally (Section~\ref{sec:eval_plan}).
    \item \textbf{Small dependency footprint.} The core is plain C with \texttt{dlopen}; optional libraries only affect the backend that needs them.
\end{itemize}

All measurements are hardware-domain measurements (package, core, DRAM, GPU) and not per-process attributions, so they include activity outside the target command.
